\section{Análisis de los Resultados}

\subsection{Comportamiento de las variables}

Las variables de resultado presentan tres características que condicionan todo el análisis. La primera es la asimetría: \texttt{cc\_max}, \texttt{cbo\_mod} y \texttt{sloc} tienen colas derechas largas, con máximos muy alejados de la mediana. La segunda es la presencia de valores extremos: un módulo puede alcanzar 533 de complejidad ciclomática o 13\,018 líneas, valores que no son errores sino expresiones de módulos excepcionalmente grandes. La tercera es la inflación de ceros en la variable independiente: 15\,567 de 15\,768 módulos no tienen hallazgos, de modo que cualquier correlación con \texttt{n\_ap} está dominada por la diferencia entre cero y uno, y no por la variación entre uno y dieciocho.

Estas características explican por qué el contraste de grupos y la correlación no cuentan la misma historia. El contraste de Mann--Whitney compara dos poblaciones y es sensible a la separación entre \emph{tener} y \emph{no tener}; la correlación de Spearman utiliza todos los valores de \texttt{n\_ap} y se atenúa cuando la mayoría se concentra en cero. Ninguno de los dos resultados es incorrecto: miden aspectos distintos de la misma asociación. El informe los presenta juntos para evitar que un único indicador sugiera una conclusión más fuerte de la que los datos permiten.

\subsection{Verificación de la pregunta de investigación}

La pregunta formulada en el anteproyecto indaga de qué manera la presencia de antipatrones arquitectónicos en aplicaciones GenAI/RAG afecta su mantenibilidad. Con el diseño ejecutado no es posible responder a la pregunta en términos de afectación causal, porque no hubo manipulación, aleatorización ni control de variables confusoras, y porque la medición es transversal. La respuesta que el estudio puede sostener es de naturaleza asociativa: en la muestra analizada, los módulos con al menos un antipatrón detectado presentan, en mediana, mayor complejidad ciclomática máxima (14 frente a 5), menor índice de mantenibilidad (39,74 frente a 68,04) y mayor acoplamiento aproximado (3 frente a 2) que los módulos sin antipatrones.

La asociación es estadísticamente significativa después de la corrección de Holm y su dirección coincide con la esperada. Su magnitud es grande para \texttt{cc\_max} y \texttt{mi} y mediana para \texttt{cbo\_mod} en el contraste de grupos, pero débil cuando se mide como correlación entre el número de hallazgos y las métricas. Además, la revisión asistida por IA sugiere que una parte no desdeñable de los hallazgos podría ser falsa, y no se dispone de validación humana independiente. Por tanto, la respuesta debe formularse con tres matices: hay asociación, la asociación es consistente en dirección y significativa, y la evidencia no permite afirmar causalidad ni cuantificar con precisión la prevalencia real.

\subsection{Verificación de los objetivos}

El primer objetivo específico consistía en formular una taxonomía conceptual de antipatrones arquitectónicos específicos para sistemas GenAI/RAG a partir de la literatura emergente y las guías de la industria. El resultado fue una taxonomía operacional de seis antipatrones (AP1--AP6), con definición, atributo afectado, regla AST, falsos positivos probables y exclusiones. La taxonomía se formuló, pero su correspondencia con la literatura quedó marcada como pendiente de verificación en varios casos, y solo cuatro de los seis antipatrones se operationalizaron. El objetivo se cumple parcialmente: hay una taxonomía utilizable y trazable, no una validación teórica completa.

El segundo objetivo específico consistía en desarrollar una herramienta de análisis estático basada en AST para automatizar la detección de la taxonomía en proyectos de GitHub. La herramienta se implementó, se ejecutó sobre 23 repositorios y 15\,769 archivos, y produjo 445 hallazgos sin errores de análisis sintáctico. El objetivo se cumple en términos de implementación y ejecución, aunque la herramienta no haya sido validada contra un estándar humano ni evaluada en cuanto a exhaustividad.

El tercer objetivo específico consistía en evaluar estadísticamente la frecuencia de aparición de los antipatrones y su correlación con métricas tradicionales de acoplamiento y complejidad sintáctica. Se calcularon descriptivos, contrastes de Mann--Whitney con tamaños de efecto, correlaciones de Spearman, terciles de tamaño y sensibilidad por repositorio. El objetivo se cumple, con las limitaciones de confusión, dependencia y comparaciones exploratorias ya señaladas. La evaluación estadística es completa en el sentido de que cubre las métricas previstas y reporta tanto significancia como magnitud, pero no es concluyente en el sentido causal.

\subsection{Verificación de las hipótesis}

La Tabla~\ref{tab:hipotesis} establece la correspondencia entre las hipótesis del anteproyecto y las decisiones estadísticas del piloto. La hipótesis de primer grado (H1) afirma que los antipatrones definidos en la taxonomía están presentes en aplicaciones open source GenAI/RAG. La regla D9 la operacionaliza exigiendo que cada antipatrón implementado aparezca en al menos dos repositorios y que al menos el 50\,\% de los repositorios tenga uno o más antipatrones. El resultado es 21 de 23 repositorios (91,3\,\%), con AP1 en 20, AP5 en 14, AP3 en 6 y AP4 en 3. Las dos condiciones se cumplen para los cuatro antipatrones implementados, por lo que D9 se registra como \emph{apoyada en el piloto}. La decisión no cubre AP2 ni AP6 y no confirma definitivamente H1.

\begin{table*}[t]
\centering
\scriptsize
\setlength{\tabcolsep}{4pt}
\caption{Correspondencia entre las hipótesis de trabajo y las decisiones estadísticas del piloto}
\label{tab:hipotesis}
\begin{tabularx}{\textwidth}{@{}c >{\raggedright\arraybackslash}p{3.4cm} c >{\raggedright\arraybackslash}X >{\raggedright\arraybackslash}p{2.2cm} >{\raggedright\arraybackslash}p{2.6cm}@{}}
\toprule
\textbf{Hip.} & \textbf{Enunciado abreviado} & \textbf{Regla} & \textbf{Evidencia principal} & \textbf{Decisión} & \textbf{Alcance} \\
\midrule
H1 & Los antipatrones de la taxonomía están presentes en aplicaciones open source GenAI/RAG. & D9 & 21 de 23 repositorios (91,3~\%) con al menos un hallazgo; AP1 en 20, AP5 en 14, AP3 en 6 y AP4 en 3 repositorios. & Apoyada en el piloto & Descriptiva; no cubre AP2 ni AP6. \\
H2 & La presencia de antipatrones se asocia con mayor deterioro de mantenibilidad y calidad estructural. & D10 & 3 de 3 métricas primarias con $p_{Holm}<0{,}001$ y magnitud $\geq$ 0,403 en la dirección esperada (\texttt{cc\_max}, \texttt{mi}, \texttt{cbo\_mod}). & Apoyada en el piloto & Asociación correlacional; no causalidad. \\
\bottomrule
\end{tabularx}
\end{table*}

La hipótesis de segundo grado (H2) afirma que la presencia de antipatrones está asociada con un mayor deterioro de las métricas de mantenibilidad y calidad estructural, particularmente en acoplamiento, complejidad ciclomática e índice de mantenibilidad. La regla D10 la operacionaliza exigiendo que al menos dos de las tres métricas primarias cumplan simultáneamente tres condiciones: $p_{Holm}<0{,}05$, dirección esperada y $|\delta|\geq0{,}147$. Las tres métricas cumplen: \texttt{cc\_max} ($p_{Holm}=4{,}85\times10^{-50}$, $\delta=+0{,}609$), \texttt{mi} ($p_{Holm}=2{,}00\times10^{-53}$, $\delta=-0{,}625$) y \texttt{cbo\_mod} ($p_{Holm}=7{,}41\times10^{-24}$, $\delta=+0{,}403$). D10 se registra como \emph{apoyada en el piloto}. La decisión es provisional y correlacional: no demuestra que los antipatrones causen el deterioro, ni que el deterioro observado sea atribuible exclusivamente a ellos.

La hipótesis de tercer grado no se planteó en el anteproyecto, y el estudio no la introduce. Su ausencia es coherente con el alcance definido: no se examinan variables como el framework, la antigüedad del repositorio o el número de colaboradores, que habrían requerido un diseño y un tamaño muestral distintos.

\subsection{Significancia estadística y magnitud del efecto}

Los valores $p$ del contraste primario son extremadamente pequeños, pero sería un error interpretarlos como evidencia de una relación fuerte. Con 15\,768 observaciones, diferencias pequeñas pueden alcanzar significancia estadística; lo que importa es el tamaño de efecto y su interpretación sustantiva. En \texttt{cc\_max} y \texttt{mi}, los $\delta$ de Cliff superan 0,6, lo que indica que la probabilidad de que un módulo con antipatrón supere a uno sin antipatrón es considerablemente mayor que la probabilidad contraria. En \texttt{cbo\_mod}, el $\delta$ de 0,403 es mediano: hay separación, pero también una superposición importante.

La correlación de Spearman matiza la lectura. Un $\rho$ de 0,119 o de $-0{,}123$ explica una fracción muy pequeña de la variación y no permite usar el número de hallazgos como predictor práctico de la métrica. La diferencia entre ambos indicadores se explica por la inflación de ceros y por la forma en que cada medida utiliza la información. El informe evita afirmar que “a mayor número de antipatrones, peor mantenibilidad” como si fuera una relación dosis-respuesta; lo que los datos muestran es una separación entre grupos y una tendencia monótona débil.

El caso de \texttt{cbo\_mod} en el tercil superior es ilustrativo. La diferencia es estadísticamente significativa ($p=0{,}00186$), pero el $\delta$ es 0,142, por debajo del umbral de 0,147, y las medianas son iguales a 3 en ambos grupos. Se trata de un resultado que la significancia aislada habría presentado como hallazgo y que el tamaño de efecto obliga a calificar como despreciable. Esta es la razón por la que el protocolo exige dirección y magnitud además del valor $p$.

\subsection{Tamaño del código, anidamiento y concentración}

El análisis por terciles de \texttt{sloc} aporta evidencia parcial de que la asociación no se explica solo por el tamaño. En el tercil medio y en el superior, \texttt{cc\_max} y \texttt{mi} mantienen diferencias significativas en la dirección esperada, con efectos pequeños o medianos. Si el tamaño fuera la única explicación, las diferencias deberían desaparecer al comparar módulos de tamaño similar; no desaparecen del todo. Sin embargo, la estratificación no es un control experimental: dentro de cada tercil sigue habiendo variación de tamaño, y otras características del repositorio, como el estilo de diseño o la madurez del proyecto, no se controlan.

El anidamiento de módulos en repositorios es una limitación más seria. Los módulos de un mismo repositorio comparten autores, dependencias, convenciones y decisiones arquitectónicas; sus métricas no son independientes. El contraste de Mann--Whitney trata cada módulo como una observación independiente, lo que puede subestimar la incertidumbre y producir valores $p$ demasiado pequeños. Un modelo multinivel con interceptos aleatorios por repositorio habría sido la alternativa metodológica adecuada, pero no se implementó. La sensibilidad agregada por repositorio intentó compensar esta limitación, pero quedó inconclusa por el reducido número de repositorios sin hallazgos.

La concentración de hallazgos en dos repositorios agrava el problema. \texttt{multimindlab/multimind-sdk} y \texttt{bojieli/ai-agent-book} aportan el 62,9\,\% de los hallazgos, de modo que los resultados globales pueden estar reflejando las prácticas de esos dos proyectos más que una propiedad general del dominio. La tasa de módulos con antipatrón de \texttt{multimindlab/multimind-sdk} (10,71\,\%) es muy superior a la de la mayoría de los repositorios, lo que sugiere un efecto de proyecto. Sin un modelo multinivel o una muestra más equilibrada, no es posible separar el efecto de módulo del efecto de repositorio.

\subsection{Alcance de la inferencia}

La evidencia permite sostener tres afirmaciones. Primera, los antipatrones implementados existen en la muestra y se distribuyen de forma desigual entre repositorios y antipatrones. Segunda, los módulos con antipatrones presentan, en conjunto, peores valores de complejidad, mantenibilidad y acoplamiento aproximado, con tamaños de efecto que van de medianos a grandes en el contraste de grupos. Tercera, la asociación persiste parcialmente al estratificar por tamaño, lo que reduce pero no elimina la plausibilidad de una explicación basada en el tamaño.

La evidencia no permite sostener cuatro afirmaciones. Primera, no permite afirmar que los antipatrones causen el deterioro. Segunda, no permite estimar la prevalencia real, porque no se conoce la precisión ni la exhaustividad de los detectores. Tercera, no permite generalizar a aplicaciones propietarias, a otros lenguajes o a repositorios que no cumplen los criterios de selección. Cuarta, no permite comparar la gravedad de los antipatrones entre sí, porque sus reglas tienen sensibilidades distintas y solo cuatro de los seis fueron implementados.

Estas distinciones no debilitan el valor del estudio; delimitan lo que un piloto empírico puede aportar. El estudio construyó un instrumento, lo ejecutó sobre datos reales, midió asociaciones con procedimientos estadísticos explícitos y documentó sus límites. Esa combinación es suficiente para justificar una investigación posterior con validación humana y diseño multinivel, pero no para cerrar la pregunta de investigación de forma definitiva.
